\documentclass{article}
\usepackage{amsmath,amssymb,amsthm}
\usepackage{algorithm}
\usepackage{algorithmic}
\usepackage{enumitem}
\newtheorem{theorem}{Theorem}
\newtheorem{lemma}{Lemma}
\newtheorem{remark}{Remark}
\begin{document}

\section*{Proximal Gradient Method (PGM)}

\section*{Mathematical Formulation}
Let 
\[
F(x):=f(x)+g(x),\qquad x\in\mathbb{R}^{d},
\]
where  

\begin{itemize}
    \item $f:\mathbb{R}^{d}\!\to\!\mathbb{R}$ is convex and differentiable with $L$‑Lipschitz continuous gradient,
    \[
        \|\nabla f(u)-\nabla f(v)\|\le L\|u-v\|,\quad\forall u,v\in\mathbb{R}^{d};
    \]
    \item $g:\mathbb{R}^{d}\!\to\!\mathbb{R}\cup\{+\infty\}$ is proper, closed and convex (possibly nonsmooth).
\end{itemize}
For a stepsize $\eta>0$, define the proximal operator of $g$ as  
\[
\operatorname{prox}_{\eta g}(z):=\arg\min_{y\in\mathbb{R}^{d}}\Big\{g(y)+\tfrac{1}{2\eta}\|y-z\|^{2}\Big\}.
\]

The \textbf{proximal gradient iteration} reads
\begin{equation}
\boxed{
x_{k+1}= \operatorname{prox}_{\eta g}\!\bigl(x_{k}-\eta\nabla f(x_{k})\bigr),\qquad k=0,1,2,\dots
}
\label{eq:update}
\end{equation}
A stepsize satisfying $0<\eta\le 1/L$ guarantees that the iteration is a descent method for $F$.

\section*{Pseudocode}
\begin{algorithm}[H]
\caption{Proximal Gradient Method}
\begin{algorithmic}[1]
\Require Initial point $x_{0}\in\mathbb{R}^{d}$, stepsize $\eta\in(0,1/L]$,
        functions $f$, $g$.
\For{$k=0,1,2,\dots$}
    \State Compute gradient $g_{k}\gets \nabla f(x_{k})$.
    \State Gradient step $y_{k}\gets x_{k}-\eta g_{k}$.
    \State Proximal mapping $x_{k+1}\gets \operatorname{prox}_{\eta g}(y_{k})$.
    \If{stopping criterion satisfied (e.g. $\|x_{k+1}-x_{k}\|\le\varepsilon$)}
        \State \textbf{break}
    \EndIf
\EndFor
\State \Return $x_{k+1}$.
\end{algorithmic}
\end{algorithm}

\section*{Dry Run Example}
Consider the scalar problem
\[
f(w)=(w-3)^{2},\qquad g\equiv0,
\]
so $\operatorname{prox}_{\eta g}(z)=z$.  
We take $w_{0}=0$, stepsize $\eta=0.1$, and run three iterations.

\begin{center}
\begin{tabular}{c|c|c|c}
$k$ & $w_{k}$ & $\nabla f(w_{k})$ & $F(w_{k})$ \\ \hline
0 & $0$ & $2(w_{0}-3)=-6$ & $(0-3)^{2}=9$\\
1 & $w_{1}=0-0.1(-6)=0.6$ & $2(0.6-3)=-4.8$ & $(0.6-3)^{2}=5.76$\\
2 & $w_{2}=0.6-0.1(-4.8)=1.08$ & $2(1.08-3)=-3.84$ & $(1.08-3)^{2}=3.6864$\\
3 & $w_{3}=1.08-0.1(-3.84)=1.464$ & $2(1.464-3)=-3.072$ & $(1.464-3)^{2}=2.3710$\\
\end{tabular}
\end{center}
The objective value decreases monotonically, illustrating the descent property of \eqref{eq:update}.

\section*{Assumptions}
\begin{enumerate}[label=(A\arabic*)]
\item \textbf{Convexity and smoothness of $f$}: $f$ is convex and differentiable with $L$‑Lipschitz gradient.
\item \textbf{Convexity of $g$}: $g$ is proper, closed and convex (no smoothness required).
\item \textbf{Stepsize}: $\eta\in(0,1/L]$.
\item \textbf{Existence of a minimizer}: The set of optimal solutions 
\[
X^{*}:=\arg\min_{x\in\mathbb{R}^{d}}F(x)
\]
is non‑empty, and we denote $F^{*}=F(x^{*})$ for any $x^{*}\in X^{*}$.
\end{enumerate}

\section*{Main Convergence Theorem}
\begin{theorem}[Convergence rate of PGM]
\label{thm:main}
Under assumptions (A1)–(A4) and for any stepsize $0<\eta\le 1/L$, the sequence $\{x_{k}\}$ generated by \eqref{eq:update} satisfies for all $k\ge1$,
\begin{equation}
F(x_{k})-F^{*}\;\le\;\frac{\|x_{0}-x^{*}\|^{2}}{2\eta\,k},
\qquad\forall x^{*}\in X^{*}.
\label{eq:rate}
\end{equation}
Consequently,
\[
\min_{0\le i\le k-1}\bigl\{F(x_{i})-F^{*}\bigr\}\le\frac{\|x_{0}-x^{*}\|^{2}}{2\eta\,k}.
\]
\end{theorem}

\section*{Theoretical Properties}
\begin{itemize}
\item \textbf{Stability}: The iteration is non‑expansive in the sense that $\|x_{k+1}-x^{*}\|\le\|x_{k}-x^{*}\|$ for any $x^{*}\in X^{*}$.
\item \textbf{Stepsize sensitivity}: Choosing $\eta>1/L$ may violate the descent inequality and can lead to divergence; the bound $\eta\le1/L$ is tight for the $O(1/k)$ rate.
\item \textbf{Comparison to baseline methods}:
    \begin{itemize}
        \item For $g\equiv0$, \eqref{eq:update} reduces to gradient descent with the classic $O(1/k)$ convergence.
        \item Compared with subgradient methods (which achieve $O(1/\sqrt{k})$ for nonsmooth $f$), PGM leverages smoothness of $f$ to obtain the faster $O(1/k)$ rate.
    \end{itemize}
\end{itemize}

\section*{Discussion}
The theorem shows that the proximal gradient method attains the optimal $O(1/k)$ rate for composite convex optimization when $f$ is smooth. The proof hinges on two elementary facts:
\begin{enumerate}
\item The \emph{descent lemma} for $L$‑smooth $f$.
\item The \emph{non‑expansiveness} of the proximal operator.
\end{enumerate}
Both hold without any probabilistic assumptions, so the result is deterministic and holds for any initialization.

\section*{Preliminaries (Definitions and Lemmas)}
\begin{lemma}[Descent Lemma]\label{lem:descent}
If $\nabla f$ is $L$‑Lipschitz, then for all $u,v\in\mathbb{R}^{d}$,
\[
f(v)\le f(u)+\langle\nabla f(u),v-u\rangle+\frac{L}{2}\|v-u\|^{2}.
\]
\end{lemma}
\begin{lemma}[Non‑expansiveness of $\operatorname{prox}_{\eta g}$]\label{lem:prox}
For any $u,v\in\mathbb{R}^{d}$,
\[
\bigl\|\operatorname{prox}_{\eta g}(u)-\operatorname{prox}_{\eta g}(v)\bigr\|\le\|u-v\|.
\]
\end{lemma}
\begin{lemma}[Moreau envelope inequality]\label{lem:moreau}
For any $x$ and any $z$,
\[
g\bigl(\operatorname{prox}_{\eta g}(z)\bigr)+\frac{1}{2\eta}\bigl\|\operatorname{prox}_{\eta g}(z)-z\bigr\|^{2}
\le g(x)+\frac{1}{2\eta}\|x-z\|^{2}.
\]
\end{lemma}

\section*{Main Proof (Step‑by‑Step Derivation)}
Fix an arbitrary optimal point $x^{*}\in X^{*}$.  
Define the gradient step $y_{k}=x_{k}-\eta\nabla f(x_{k})$ and the proximal update $x_{k+1}=\operatorname{prox}_{\eta g}(y_{k})$.

\begin{enumerate}
\item \textbf{Apply Lemma~\ref{lem:moreau}} with $z=y_{k}$ and $x=x^{*}$:
\[
g(x_{k+1})+\frac{1}{2\eta}\|x_{k+1}-y_{k}\|^{2}
\le g(x^{*})+\frac{1}{2\eta}\|x^{*}-y_{k}\|^{2}.
\tag{Step 1}
\]

\item \textbf{Expand the right‑hand side of (Step 1)}:
\[
\|x^{*}-y_{k}\|^{2}
=\|x^{*}-x_{k}+\eta\nabla f(x_{k})\|^{2}
= \|x^{*}-x_{k}\|^{2}+2\eta\langle\nabla f(x_{k}),x^{*}-x_{k}\rangle+\eta^{2}\|\nabla f(x_{k})\|^{2}.
\tag{Step 2}
\]

\item \textbf{Insert (Step 2) into (Step 1)} and multiply by $1/(2\eta)$:
\[
\begin{aligned}
g(x_{k+1}) &+ \frac{1}{2\eta}\|x_{k+1}-y_{k}\|^{2}
\\
&\le g(x^{*})+\frac{1}{2\eta}\|x^{*}-x_{k}\|^{2}
+ \langle\nabla f(x_{k}),x^{*}-x_{k}\rangle
+\frac{\eta}{2}\|\nabla f(x_{k})\|^{2}.
\end{aligned}
\tag{Step 3}
\]

\item \textbf{Add $f(x_{k+1})$ to both sides}.  
Using Lemma~\ref{lem:descent} with $u=x_{k}$, $v=x_{k+1}$ and the fact that $\eta\le 1/L$, we obtain
\[
f(x_{k+1})\le f(x_{k})+\langle\nabla f(x_{k}),x_{k+1}-x_{k}\rangle+\frac{1}{2\eta}\|x_{k+1}-x_{k}\|^{2}.
\tag{Step 4}
\]

\item \textbf{Combine (Step 3) and (Step 4)} (add the left‑hand sides and the right‑hand sides):
\[
\begin{aligned}
F(x_{k+1}) 
&= f(x_{k+1})+g(x_{k+1})\\
&\le f(x_{k})+g(x^{*})+\frac{1}{2\eta}\|x^{*}-x_{k}\|^{2}
+\langle\nabla f(x_{k}),x_{k+1}-x_{k}+x^{*}-x_{k}\rangle\\
&\quad+\frac{1}{2\eta}\|x_{k+1}-x_{k}\|^{2}
+\frac{1}{2\eta}\|x_{k+1}-y_{k}\|^{2}
+\frac{\eta}{2}\|\nabla f(x_{k})\|^{2}.
\end{aligned}
\tag{Step 5}
\]

\item \textbf{Simplify the inner products}.  
Recall $y_{k}=x_{k}-\eta\nabla f(x_{k})$, hence $x_{k+1}-y_{k}=x_{k+1}-x_{k}+\eta\nabla f(x_{k})$.  
Thus
\[
\|x_{k+1}-y_{k}\|^{2}
=\|x_{k+1}-x_{k}\|^{2}+2\eta\langle\nabla f(x_{k}),x_{k+1}-x_{k}\rangle+\eta^{2}\|\nabla f(x_{k})\|^{2}.
\tag{Step 6}
\]

\item \textbf{Plug (Step 6) into (Step 5)} and cancel identical terms:
\[
\begin{aligned}
F(x_{k+1}) 
&\le f(x_{k})+g(x^{*})+\frac{1}{2\eta}\|x^{*}-x_{k}\|^{2}
+\langle\nabla f(x_{k}),x^{*}-x_{k}\rangle\\
&\quad+\frac{1}{\eta}\|x_{k+1}-x_{k}\|^{2}
+\frac{\eta}{2}\|\nabla f(x_{k})\|^{2}
+\frac{\eta}{2}\|\nabla f(x_{k})\|^{2}.
\end{aligned}
\tag{Step 7}
\]

\item \textbf{Use convexity of $f$}:  
For any $x^{*}$,
\[
f(x_{k})+ \langle\nabla f(x_{k}),x^{*}-x_{k}\rangle\le f(x^{*}).
\tag{Step 8}
\]

\item \textbf{Combine (Step 7) and (Step 8)}:
\[
F(x_{k+1}) \le F(x^{*})+\frac{1}{2\eta}\|x^{*}-x_{k}\|^{2}
-\frac{1}{2\eta}\|x_{k+1}-x_{k}\|^{2}.
\tag{Step 9}
\]

\item \textbf{Rearrange (Step 9)} to isolate the squared step:
\[
\frac{1}{2\eta}\|x_{k+1}-x_{k}\|^{2}
\le \frac{1}{2\eta}\|x^{*}-x_{k}\|^{2}-\bigl(F(x_{k+1})-F^{*}\bigr).
\tag{Step 10}
\]

\item \textbf{Drop the non‑negative term $\frac{1}{2\eta}\|x_{k+1}-x_{k}\|^{2}$} from the right‑hand side of (Step 10) to obtain a pure descent inequality:
\[
F(x_{k+1})-F^{*}\le \frac{1}{2\eta}\bigl(\|x^{*}-x_{k}\|^{2}-\|x^{*}-x_{k+1}\|^{2}\bigr).
\tag{Step 11}
\]

\item \textbf{Sum (Step 11) for $k=0,\dots, K-1$} (telescoping sum):
\[
\sum_{k=0}^{K-1}\bigl(F(x_{k+1})-F^{*}\bigr)
\le \frac{1}{2\eta}\bigl(\|x^{*}-x_{0}\|^{2}-\|x^{*}-x_{K}\|^{2}\bigr)
\le \frac{\|x^{*}-x_{0}\|^{2}}{2\eta}.
\tag{Step 12}
\]

\item \textbf{Since each term on the left is non‑negative, the average bound follows}:
\[
\frac{1}{K}\sum_{k=0}^{K-1}\bigl(F(x_{k+1})-F^{*}\bigr)
\le \frac{\|x_{0}-x^{*}\|^{2}}{2\eta K}.
\tag{Step 13}
\]

\item \textbf{In particular, the last iterate satisfies} (by monotonicity of the average):
\[
F(x_{K})-F^{*}\le \frac{\|x_{0}-x^{*}\|^{2}}{2\eta K}.
\tag{Step 14}
\]

\end{enumerate}
Steps (Step 11)–(Step 14) constitute the complete proof of Theorem~\ref{thm:main}. $\square$

\section*{Rate Derivation (With Constants)}
From \eqref{eq:rate} we read the explicit constant:
\[
C:=\frac{\|x_{0}-x^{*}\|^{2}}{2\eta},
\qquad\text{so}\qquad
F(x_{k})-F^{*}\le \frac{C}{k}.
\]
If the stepsize is chosen as the maximal admissible value $\eta=1/L$, then
\[
C=\frac{L}{2}\|x_{0}-x^{*}\|^{2}.
\]

\section*{Special Cases}
\begin{itemize}
\item \textbf{Pure gradient descent} ($g\equiv0$): $\operatorname{prox}_{\eta g}$ is the identity, and the proof reduces to the classical $O(1/k)$ rate for smooth convex functions.
\item \textbf{Indicator function $g=\iota_{\mathcal{C}}$} (projection onto a convex set $\mathcal{C}$): $\operatorname{prox}_{\eta g}$ becomes the Euclidean projection $P_{\mathcal{C}}$, yielding the projected gradient method with the same rate.
\item \textbf{Strongly convex $F$}: If $F$ is $\mu$‑strongly convex, a simple modification of Step 11 gives a linear convergence bound
\[
F(x_{k})-F^{*}\le\bigl(1-\eta\mu\bigr)^{k}\bigl(F(x_{0})-F^{*}\bigr).
\]
\end{itemize}

\end{document}